\chapter{PENDAHULUAN}
\section{Latar Belakang}
Uraikan konteks masalah, kesenjangan penelitian, dan alasan ilmiah penelitian perlu dilakukan. Dukungan literatur yang relevan wajib disitasi dengan perintah \texttt{\textbackslash cite} menggunakan kunci pustaka pada berkas \texttt{Referensi.bib}.
\section{Rumusan Masalah}
Nyatakan pertanyaan penelitian yang spesifik, terukur, dan dapat dijawab melalui metode yang dipilih.
\section{Tujuan Penelitian}
Nyatakan tujuan yang berkorespondensi langsung dengan rumusan masalah.
\section{Manfaat Penelitian}
Jelaskan kontribusi teoretis atau penerapan yang diharapkan.
\section{Batasan Masalah}
Tegaskan ruang lingkup, asumsi, dan keterbatasan penelitian.
